\documentclass[11pt]{article}
\usepackage{amsmath,amssymb,amsthm,mathtools}
\usepackage[margin=1in]{geometry}
\theoremstyle{plain}
\newtheorem{theorem}{Theorem}
\newtheorem{lemma}{Lemma}
\theoremstyle{remark}
\newtheorem{remark}{Remark}
\newcommand{\E}{\mathbb{E}}\newcommand{\PP}{\mathbb{P}}
\newcommand{\To}{\Rightarrow}

\title{Formalization brief: the martingale central limit theorem\\
(discrete core of Ethier--Kurtz Thm.~7.1.4)}
\date{5 July 2026}

\begin{document}
\maketitle

\noindent\textbf{Goal.} A new self-contained file (suggest
\texttt{TypeDDecouplingMartingaleCLT.lean}) proving a martingale central limit
theorem usable as the probabilistic input behind \texttt{prop\_twophase}
(\texttt{TypeDDecouplingCrossover.lean}). The paper cites Ethier--Kurtz
Theorem~7.1.4 (martingale functional CLT). The \underline{full} path-space statement
(Skorokhod topology, c\`adl\`ag martingales) is out of scope: Mathlib has no
Skorokhod space. The target is the \underline{fixed-time, finite-dimensional core},
which is what the paper's application actually evaluates. Whole project must
build; standard axioms only; do not modify existing files except (optionally) a
docstring cross-reference in the Crossover file.

\section*{Tier 1 (primary goal): McLeish's martingale CLT}

\begin{theorem}[martingale difference array CLT; McLeish 1974]\label{thm:mcleish}
For each $n$ let $(X_{n,j})_{j=1}^{k_n}$ be a square-integrable martingale
difference array with respect to filtrations
$(\mathcal F_{n,j})_j$: $\E[X_{n,j}\mid\mathcal F_{n,j-1}]=0$. Assume:
\begin{itemize}
\item[(a)] \textup{(vanishing maximal jump)}
  $\max_{j}|X_{n,j}|\to0$ in probability, and
  $\sup_n\E\big[\max_j X_{n,j}^2\big]<\infty$;
\item[(b)] \textup{(bracket convergence)}
  $\sum_{j=1}^{k_n}X_{n,j}^2\ \to\ \sigma^2$ in probability, $\sigma^2\ge0$
  a deterministic constant.
\end{itemize}
Then $S_n:=\sum_{j=1}^{k_n}X_{n,j}\To N(0,\sigma^2)$ in distribution.
\end{theorem}

\noindent\textbf{Proof route (McLeish's product trick; recommended).} Set
$T_n(u):=\prod_j\big(1+\mathrm iuX_{n,j}\big)$. Then:
(1) $\E[T_n(u)]=1$ exactly, by the martingale property, expanding the product
in the natural filtration order.
(2) $\log(1+\mathrm ix)=\mathrm ix+\tfrac{x^2}2 r(x)$ with $|r|\le1$ bounded and
$r(x)\to-1$ as $x\to0$, so
$T_n(u)=\exp\big(\mathrm iuS_n+\tfrac{u^2}2\sum_jX_{n,j}^2\,r_j\big)$ where the
correction converges: by (a)+(b),
$\exp(\mathrm iuS_n)\,e^{u^2\sigma^2/2}-T_n(u)\to0$ in $L^1$
(uniform integrability of $T_n(u)$ from
$|T_n(u)|^2=\prod(1+u^2X_{n,j}^2)\le\exp(u^2\sum X_{n,j}^2)$ after a standard
truncation of $\sum X_{n,j}^2$ at a constant, which is harmless by (b)).
(3) Hence $\E[e^{\mathrm iuS_n}]\to e^{-u^2\sigma^2/2}$ for every $u$, and
L\'evy continuity concludes. Mathlib's characteristic-function API
(\texttt{charFun}, L\'evy's continuity theorem) is available; if the exact
truncation lemma is painful, introducing the stopped array
$\tilde X_{n,j}:=X_{n,j}\mathbf 1\{\sum_{i<j}X_{n,i}^2\le2\sigma^2\}$ (still a
martingale difference array, equal to $X$ with probability $\to1$) is the
standard device.

\section*{Tier 2: the bivariate version the paper uses}

\begin{theorem}[joint CLT with diagonal brackets]\label{thm:joint}
Let $(X_{n,j}),(Y_{n,j})$ be martingale difference arrays with respect to a
common filtration, each satisfying (a), with
$\sum_jX_{n,j}^2\to\sigma_X^2$, $\sum_jY_{n,j}^2\to\sigma_Y^2$ and
$\sum_jX_{n,j}Y_{n,j}\to0$, all in probability. Then
$(\sum_jX_{n,j},\ \sum_jY_{n,j})\To N(0,\sigma_X^2)\otimes N(0,\sigma_Y^2)$
(independent Gaussians).
\end{theorem}

\noindent Route: Cram\'er--Wold through Tier~1 --- for fixed $(a,b)$ the array
$aX_{n,j}+bY_{n,j}$ satisfies (a),(b) with limit
$a^2\sigma_X^2+b^2\sigma_Y^2$, so every linear combination is asymptotically
normal with the product covariance; conclude by characteristic functions of the
pair (two-dimensional L\'evy continuity, or directly: the 2d charFun evaluated
at $(u,v)$ is the 1d charFun of the combined array at $1$).

This is exactly the shape used by \texttt{prop:twophase}: the compensated jump
martingales $(\bar S,\bar R)$ evaluated at time $T-\tau$, with diagonal brackets
$2(1+q^2)(T-\tau)+O_\PP(\sqrt T)$ and cross bracket $\eps\Lambda_T=o_\PP(T)$,
normalized by $\sqrt{2T}$.

\section*{Tier 3 (stretch; only if Tiers 1--2 complete): continuous-time corollary}

For a pure-jump square-integrable martingale $M$ with bounded jumps
$|\Delta M|\le b$ and finitely many jumps on $[0,T]$ a.s.\ (the paper's setting:
compensated finite-rate Markov jump processes), the discretization
$X_{n,j}:=M(jT/n)-M((j-1)T/n)$ is a martingale difference array; conditions
(a),(b) for it follow from $\langle M\rangle_T/c_T\to\sigma^2$ and
$b/\sqrt{c_T}\to0$ under the normalization $M(T)/\sqrt{c_T}$. State and prove
this reduction if the process encoding permits; otherwise state it as a named
hypothesis bridging to Tier~2 and record the interface in a docstring. Do not
attempt Skorokhod-space path convergence.

\begin{remark}[interface and fallbacks]
(1) Do not modify \texttt{prop\_twophase} or its \texttt{sorry}; if Tier~2
completes, add only a docstring in the Crossover file noting that the
probabilistic core (the joint martingale CLT) is now
\texttt{thm\_joint} in the new file, with the time discretization and the
pre/post-split decomposition remaining in the leaf.
(2) Sanctioned fallbacks, in order: truncation lemma as a separate named
lemma; conditions (a) strengthened to a deterministic bound
$\max_j|X_{n,j}|\le b_n\to0$ (sufficient for the paper: jumps are bounded by
$2/\sqrt{2T}$ after normalization); if L\'evy continuity in Mathlib lacks a
needed form, state that single step as a named hypothesis and prove everything
else.
(3) Constants and generality are free: real-valued arrays suffice; no need for
vector-valued beyond the bivariate Tier~2; $\sigma^2=0$ (degenerate limit) may
be excluded by hypothesis $\sigma^2>0$ if it simplifies, since the paper's
brackets are nondegenerate.
\end{remark}

\end{document}
